\documentclass[12pt, twoside]{article}
%\documentclass[12pt, twoside]{article}
\usepackage[letterpaper, margin=1in, headsep=0.2in]{geometry}
\setlength{\headheight}{0.6in}
%\usepackage[english]{babel}
\usepackage[utf8]{inputenc}
\usepackage{microtype}
\usepackage{amsmath}
\usepackage{amssymb}
%\usepackage{amsfonts}
\usepackage[nomessages]{fp} %\FPeval{\var-name}{2*sin(pi/6)}
\usepackage{siunitx} %units in math. eg 20\milli\meter
\usepackage{yhmath} % for arcs, overparenth command
\usepackage{tikz} %graphics
\usetikzlibrary{quotes, angles, arrows, arrows.meta}
\usepackage{pgfplots}
\usepgfplotslibrary{fillbetween}
\pgfplotsset{compat=1.16}
\usepackage{graphicx} %consider setting \graphicspath{{images/}}
\usepackage{parskip} %no paragraph indent
\usepackage{enumitem}
\usepackage{multicol}
\usepackage{venndiagram}

\usepackage{fancyhdr}
\pagestyle{fancy}
\fancyhf{}
\renewcommand{\headrulewidth}{0pt} % disable the underline of the header
\raggedbottom
\hfuzz=2mm %suppresses overfull box warnings

\usepackage{hyperref}

\fancyhead[LE]{\thepage}
\fancyhead[RO]{Markscheme}
\fancyhead[LO]{La Scuola d'Italia / Dr.\ Huson / IB Math 11 \\* 26 October 2026}

\begin{document}
\subsubsection*{2.4 Pre-Quiz: Lines, Exponents and Radicals --- Markscheme}

\emph{Total: 77 marks.}

\begin{enumerate}[itemsep=0.3cm]

%--- Problem 1: Triangle ABC: gradient of AB, perpendicular from C, foot D, height CD in surd form, area [14 marks] ---
\item[1.] [Maximum mark: 14]

\textbf{(a)} $m_{AB} = \dfrac{5 - 1}{6 - (-2)}$ \hfill \textbf{M1}

$m_{AB} = \dfrac{4}{8} = \dfrac{1}{2}$ \hfill \textbf{A1}

\emph{A bare answer earns M0A1.}

\textbf{(b)} $m_{L} = -2$ \hfill \textbf{A1}

$y - (-2) = -2(x - 7)$ \hfill \textbf{M1}

$y = -2x + 12$ \hfill \textbf{A1}

\emph{A1 for the negative reciprocal of their gradient. M1 for substituting
their gradient and $C(7, -2)$ into any form of the equation of a line. The
final A1 needs the form $y = mx + c$. FT from (a).}

\emph{Wrong method: the perpendicular bisector of $AB$,
$y = -2x + 7$, passes through the midpoint of $AB$, not through $C$: A1M0A0.}

\textbf{(c)} line $AB$: $y - 1 = \frac{1}{2}(x + 2)$, so $y = \frac{1}{2}x + 2$ \hfill \textbf{A1}

$-2x + 12 = \frac{1}{2}x + 2$ \hfill \textbf{M1}

$\frac{5}{2}x = 10$

$x = 4$, \quad $y = 4$ \qquad $D(4,\ 4)$ \hfill \textbf{A1A1}

\emph{M1 for equating their two lines, or for any valid method of solving
them together (substitution, elimination, GDC). A1 for each coordinate. FT
from (b). Wrong method: the midpoint of $AB$, $(2,\ 3)$, is not the foot of
the perpendicular from $C$: A1 for the equation of $AB$ if seen, then
M0A0A0.}

\textbf{(d)} $CD = \sqrt{(7 - 4)^{2} + (-2 - 4)^{2}} = \sqrt{45}$ \hfill \textbf{M1}

$CD = 3\sqrt{5}$ \hfill \textbf{A1}

\emph{Accept $a = 3$. $\sqrt{45}$ or $6.71$ alone earns A0. FT from (c).
Wrong method: $CA = 3\sqrt{10}$ or $CB = 5\sqrt{2}$ is not the
perpendicular distance: M0A0.}

\textbf{(e)} $AB = \sqrt{8^{2} + 4^{2}} = \sqrt{80} = 4\sqrt{5}$ \hfill \textbf{A1}

$\text{Area} = \frac{1}{2}\left(4\sqrt{5}\right)\left(3\sqrt{5}\right)$ \hfill \textbf{M1}

$\text{Area} = 30$ \hfill \textbf{A1}

\emph{M1 for $\frac{1}{2} \times \text{base} \times \text{height}$ with
$AB$ as the base and $CD$ as the height. Accept any other correct method,
e.g.\ enclosing the triangle in a rectangle and subtracting. FT from (d).}

\newpage

%--- Problem 2: Car hire: linear cost model from two data points, interpret intercept, break-even, compare at 9 days [10 marks] ---
\item[2.] [Maximum mark: 10]

\textbf{(a)} $m = \dfrac{305 - 165}{7 - 3}$ \hfill \textbf{M1}

$m = 35$ \hfill \textbf{A1}

\textbf{(b)} $165 = 35(3) + c$ \hfill \textbf{M1}

$c = 60$ \hfill \textbf{A1}

\emph{Accept solving the system $3m + c = 165$, $7m + c = 305$ for both
values at once. FT from (a).}

\textbf{(c)} $c = 60$ is the fixed fee: company $A$ charges $60$ euros
before any days are added. \hfill \textbf{A1}

\emph{Accept ``the cost for $0$ days'' or ``the fixed charge''. Do not accept
``the $y$-intercept'' with no reference to the context.}

\textbf{(d)} $35d + 60 = 45d + 10$ \hfill \textbf{M1}

$10d = 50$

$d = 5$ \hfill \textbf{A1}

\emph{Accept a graphical or table method on the GDC. FT from (a) and (b).}

\textbf{(e)} $C_{A}(9) = 35(9) + 60 = 375$, \quad $C_{B}(9) = 45(9) + 10 = 415$ \hfill \textbf{M1}

company $A$ is cheaper \hfill \textbf{A1}

she saves $415 - 375 = 40$ euros \hfill \textbf{A1}

\emph{M1 for evaluating both costs at $d = 9$. FT from (a) and (b).}

\newpage

%--- Problem 3: Solve a 2x2 system; lines with parameter k: parallel (no solution), perpendicular, intersection [13 marks] ---
\item[3.] [Maximum mark: 13]

\textbf{(a)} $3x - 4y = 18$ and $10x + 4y = 8$, adding: $13x = 26$ \hfill \textbf{(M1)}

$x = 2$ \hfill \textbf{A1}

$y = -3$ \hfill \textbf{A1}

\emph{M1 for any valid method: elimination, substitution, or the GDC
simultaneous-equation solver. A correct answer from the GDC with no other
working earns full marks.}

\textbf{(b)} $L_{2}$: $y = \frac{1}{2}x - \frac{5}{6}$, gradient $\frac{1}{2}$ \hfill \textbf{A1}

$L_{1}$: $y = -\frac{2}{k}x + \frac{6}{k}$, so $-\frac{2}{k} = \frac{1}{2}$ \hfill \textbf{M1}

$k = -4$ \hfill \textbf{A1}

\emph{M1 for equating the gradients of the two lines. Accept comparing
coefficients: $\frac{2}{3} = \frac{k}{-6}$.}

\textbf{(c)} no solutions \hfill \textbf{A1}

the lines are parallel and are not the same line ($y$-intercepts
$-\frac{3}{2}$ and $-\frac{5}{6}$), so they never meet \hfill \textbf{R1}

\emph{Accept ``parallel lines do not intersect''. Do not award A1R0.}

\textbf{(d)} $\left(-\frac{2}{k}\right) \times \frac{1}{2} = -1$ \hfill \textbf{M1}

$k = 1$ \hfill \textbf{A1}

\emph{M1 for using the product of the gradients $= -1$, or the gradient of
$L_{1}$ equal to $-2$. FT from the gradient of $L_{2}$ in (b).}

\textbf{(e)} $2x + y = 6$, \quad $3x - 6y = 5$

$3x - 6(6 - 2x) = 5$ \hfill \textbf{M1}

$15x = 41$, \quad $x = \dfrac{41}{15}$ \hfill \textbf{A1}

$y = 6 - 2\left(\dfrac{41}{15}\right) = \dfrac{8}{15}$ \hfill \textbf{A1}

\emph{M1 for any valid method, including the GDC solver. The question asks
for exact values: decimal coordinates $(2.73,\ 0.533)$ earn A1A0. FT from
(d).}

\newpage

%--- Problem 4: Exponent laws with negative/fractional powers; same-base equation; hidden quadratic in 2^x; GDC solution of 5^x = 40 [15 marks] ---
\item[4.] [Maximum mark: 15]

\textbf{(a)} $\left(2x^{3}y^{-2}\right)^{3} = 8x^{9}y^{-6}$ \hfill \textbf{M1}

$\dfrac{8x^{9}y^{-6}}{4x^{5}y^{-4}} = 2x^{4}y^{-2}$ \hfill \textbf{A1}

$= \dfrac{2x^{4}}{y^{2}}$ \hfill \textbf{A1}

\emph{M1 for raising every factor of the bracket to the power $3$. The final
A1 needs positive exponents only.}

\textbf{(b)} $27^{-\frac{2}{3}} = \dfrac{1}{\left(\sqrt[3]{27}\right)^{2}} = \dfrac{1}{9}$ \hfill \textbf{A1}

$\left(x^{6}\right)^{-\frac{2}{3}} = x^{-4}$, \quad so the expression is $\dfrac{1}{9}x^{-4}$ \hfill \textbf{A1}

\emph{Accept $a = \frac{1}{9}$, $n = -4$. Do not accept $0.111x^{-4}$:
$a$ must be exact.}

\textbf{(c)} $\left(3^{2}\right)^{x + 1} = \left(3^{-3}\right)^{x}$ \hfill \textbf{M1}

$2x + 2 = -3x$ \hfill \textbf{A1}

$x = -\dfrac{2}{5}$ \hfill \textbf{A1}

\emph{M1 for writing both sides as powers of $3$. Accept $x = -0.4$. A GDC
answer with no working earns M0A0A1.}

\textbf{(d)} $4^{x} = \left(2^{2}\right)^{x} = \left(2^{x}\right)^{2} = u^{2}$ \hfill \textbf{A1}

so $2u^{2} - 17u + 8 = 0$ \hfill \textbf{AG}

\textbf{(e)} $(2u - 1)(u - 8) = 0$ \hfill \textbf{M1}

$u = \frac{1}{2}$ or $u = 8$ \hfill \textbf{A1}

$2^{x} = \frac{1}{2} \Rightarrow x = -1$ \hfill \textbf{A1}

$2^{x} = 8 \Rightarrow x = 3$ \hfill \textbf{A1}

\emph{M1 for solving the quadratic in $u$ by any method (factorizing,
formula, GDC). Stopping at $u = \frac{1}{2}, 8$ earns M1A1A0A0. Reading
$2^{x}$ as $2x$ gives $x = \frac{1}{4}, 4$: A0A0 for those values. ``Hence'' requires
the substitution: reading both roots from a GDC graph of the original
equation, with no work in $u$, earns M0A0A1A1.}

\textbf{(f)} graph $y = 5^{x}$ and $y = 40$ and find the intersection \hfill \textbf{(M1)}

$x = 2.29$ \quad ($2.29203\ldots$) \hfill \textbf{A1}

\emph{The M1 is implied by a correct answer, so a bare $2.29$ earns 2/2. A1 for $2.29$ or the display
with an ellipsis; $2.292$ with no ellipsis earns A0.}

\newpage

%--- Problem 5: Surds: collect like surds, square a binomial, rationalize with conjugate, linear equation with surd coefficients, rectangle width [13 marks] ---
\item[5.] [Maximum mark: 13]

\textbf{(a)} $5\sqrt{3} + 4\sqrt{3} - 6\sqrt{3}$ \hfill \textbf{M1}

$= 3\sqrt{3}$ \hfill \textbf{A1}

\emph{M1 for writing at least two of the three surds as multiples of
$\sqrt{3}$.}

\textbf{(b)} $\left(2\sqrt{5}\right)^{2} - 2\left(2\sqrt{5}\right)(3) + 3^{2} = 20 - 12\sqrt{5} + 9$ \hfill \textbf{M1}

$= 29 - 12\sqrt{5}$ \hfill \textbf{A1}

\emph{M1 for an expansion with all three (or four) terms.}

\textbf{(c)} $\dfrac{\sqrt{3} + 1}{\sqrt{3} - 1} \times \dfrac{\sqrt{3} + 1}{\sqrt{3} + 1}$ \hfill \textbf{M1}

$= \dfrac{3 + 2\sqrt{3} + 1}{3 - 1} = \dfrac{4 + 2\sqrt{3}}{2}$ \hfill \textbf{A1}

$= 2 + \sqrt{3}$ \hfill \textbf{A1}

\emph{M1 for multiplying numerator and denominator by the conjugate
$\sqrt{3} + 1$. A1 for a correct numerator and the denominator $2$. Accept
$a = 2$, $b = 1$.}

\textbf{(d)} $x\sqrt{8} - x\sqrt{2} = 8$ \hfill \textbf{M1}

$2x\sqrt{2} - x\sqrt{2} = 8$, \quad so $x\sqrt{2} = 8$ \hfill \textbf{A1}

$x = \dfrac{8}{\sqrt{2}} = \dfrac{8\sqrt{2}}{2} = 4\sqrt{2}$ \hfill \textbf{A1}

\emph{M1 for collecting the $x$-terms on one side. A1 for
$x\sqrt{2} = 8$, or $x\left(\sqrt{8} - \sqrt{2}\right) = 8$ with
$\sqrt{8} = 2\sqrt{2}$ seen. The final A1 needs the form $a\sqrt{2}$;
$\frac{8}{\sqrt{2}}$ or $5.66$ earns A0.}

\textbf{(e)} $\text{width} = \dfrac{7 + 7\sqrt{2}}{3 + \sqrt{2}} \times \dfrac{3 - \sqrt{2}}{3 - \sqrt{2}}$ \hfill \textbf{M1}

$= \dfrac{21 - 7\sqrt{2} + 21\sqrt{2} - 14}{9 - 2} = \dfrac{7 + 14\sqrt{2}}{7}$ \hfill \textbf{A1}

$= 1 + 2\sqrt{2}$ \hfill \textbf{A1}

\emph{Accept the system method: $(3 + \sqrt{2})(a + b\sqrt{2}) = 7 + 7\sqrt{2}$
gives $3a + 2b = 7$ and $a + 3b = 7$ (M1A1), so $a = 1$, $b = 2$ (A1).}

\newpage

%--- Problem 6: Line through surd points: gradient sqrt2, equation, exact x-intercept, intersection with y = 2x (rationalize), length AB [12 marks] ---
\item[6.] [Maximum mark: 12]

\textbf{(a)} $m = \dfrac{7 - 3}{3\sqrt{2} - \sqrt{2}}$ \hfill \textbf{M1}

$m = \dfrac{4}{2\sqrt{2}} = \dfrac{2}{\sqrt{2}} = \dfrac{2\sqrt{2}}{2}$ \hfill \textbf{A1}

$m = \sqrt{2}$ \hfill \textbf{AG}

\emph{A1 for a correct step that removes the surd from the denominator, or
for $\frac{2}{\sqrt{2}} = \frac{\sqrt{2}\sqrt{2}}{\sqrt{2}}$. A decimal
check ($1.414\ldots$) does not show the result: A0.}

\textbf{(b)} $y - 3 = \sqrt{2}\left(x - \sqrt{2}\right)$ \hfill \textbf{M1}

$y = \sqrt{2}x - 2 + 3 = \sqrt{2}x + 1$ \hfill \textbf{A1}

\emph{M1 for substituting $\sqrt{2}$ and either point into any form of the
equation of a line. The A1 needs $\sqrt{2} \times \sqrt{2} = 2$ and the
form $y = mx + c$.}

\textbf{(c)} $0 = \sqrt{2}x + 1$ \hfill \textbf{M1}

$x = -\dfrac{1}{\sqrt{2}} = -\dfrac{\sqrt{2}}{2}$ \hfill \textbf{A1}

\emph{Accept $\left(-\frac{\sqrt{2}}{2},\ 0\right)$. $-\frac{1}{\sqrt{2}}$
or $-0.707$ earns A0. FT from (b).}

\textbf{(d)} $\sqrt{2}x + 1 = 2x$, \quad $x\left(2 - \sqrt{2}\right) = 1$ \hfill \textbf{M1}

$x = \dfrac{1}{2 - \sqrt{2}} \times \dfrac{2 + \sqrt{2}}{2 + \sqrt{2}} = \dfrac{2 + \sqrt{2}}{4 - 2}$ \hfill \textbf{M1}

$x = 1 + \dfrac{1}{2}\sqrt{2}$ \hfill \textbf{A1}

$y = 2x = 2 + \sqrt{2}$ \hfill \textbf{A1}

\emph{First M1 for equating the two lines and collecting the $x$-terms.
Second M1 for multiplying by the conjugate $2 + \sqrt{2}$. Accept
$x = \frac{2 + \sqrt{2}}{2}$. Decimal coordinates
$(1.71,\ 3.41)$ earn M1M0A0A0. FT from (b).}

\textbf{(e)} $AB = \sqrt{\left(3\sqrt{2} - \sqrt{2}\right)^{2} + (7 - 3)^{2}} = \sqrt{8 + 16}$ \hfill \textbf{M1}

$AB = \sqrt{24} = 2\sqrt{6}$ \hfill \textbf{A1}

\emph{The A1 needs simplest surd form: $\sqrt{24}$ or $4.90$ earns A0.}

\end{enumerate}
\end{document}
